\def\uadectitulo{Anatomía humana}
\documentclass[12pt,letterpaper,oneside]{article}
\usepackage[T1]{fontenc}
\usepackage[utf8]{inputenc}
\usepackage[spanish,es-nodecimaldot]{babel}
\usepackage{lmodern}
\usepackage[margin=2.5cm]{geometry}
\usepackage{amsmath,amssymb,graphicx,booktabs,longtable,array,enumitem}
\usepackage[most]{tcolorbox}
\usepackage{tikz}
\usetikzlibrary{babel,arrows.meta,positioning}
\usepackage{setspace}
\usepackage[hidelinks]{hyperref}
\onehalfspacing
\setlength{\emergencystretch}{3em}
\providecommand{\uadectitulo}{Material de estudio}
\providecommand{\uadecautor}{Material de estudio AulaTeX}
\newcommand{\portadaUAdeC}{%
  \hypersetup{pageanchor=false}
  \begin{titlepage}
  \centering
  {\Large Universidad Autónoma de Coahuila\par}
  \vspace{0.7cm}
  {\large Preparación para ingreso a Médico Cirujano\par}
  \vspace{1.8cm}
  {\LARGE\bfseries\uadectitulo\par}
  \vspace{1cm}
  {\large Actividad 1: síntesis y representación conceptual\par}
  {\large Actividad 2: práctica y soluciones justificadas\par}
  \vfill
  {\uadecautor\par}
  {Revisión: 10 de octubre de 2026\par}
  \vspace{0.8cm}
  \begin{minipage}{0.87\linewidth}
  \small Material formativo local. No constituye una guía, plantilla o examen oficial de UAdeC o CENEVAL. La cobertura y dificultad no están certificadas frente a una convocatoria de ingreso. No sustituye orientación ni atención clínica.
  \end{minipage}
  \end{titlepage}
  \hypersetup{pageanchor=true}
  \tableofcontents\clearpage
  \section*{Alcance y procedencia}
  Las actividades siguen las consignas locales conservadas en esta materia. Los reactivos son de práctica y las respuestas son propuestas razonadas, no resultados de un intento del aspirante. Las fuentes de apoyo se identifican al final. Los PDF bibliográficos locales se inventarían por archivo y hash; su disponibilidad no acredita lectura íntegra ni permiso de redistribución.
}
\begin{document}
\portadaUAdeC
\section{Actividad 1: síntesis de investigación}
\subsection{Organización estructural y posición}
La organización se estudia desde el nivel químico y celular hasta tejidos, órganos, sistemas y organismo. La posición anatómica permite describir relaciones sin depender de la postura momentánea. Sagital divide en porciones derecha e izquierda; solo el plano medio las divide en mitades iguales. Proximal y distal expresan distancia al origen o punto de unión y no se deben restringir de manera absoluta a extremidades.
\subsection{Cavidades y regiones}
El mediastino pertenece al tórax y contiene el corazón rodeado por el pericardio. La cavidad pericárdica es el espacio potencial entre hojas serosas, no un recipiente hueco ocupado literalmente por el corazón. Los pulmones están rodeados por las pleuras, a los lados del mediastino. Para localizar regiones abdominales se usan cuatro cuadrantes o nueve regiones; no implican compartimentos separados por paredes.
\subsection{Clasificación y función}
Los huesos largos, cortos, planos, irregulares y sesamoideos se clasifican por forma. El sistema óseo sostiene, protege, permite movimiento, almacena minerales y participa en hematopoyesis. El esqueleto axial y el apendicular son divisiones funcionales, no tipos de tejido. La tibia soporta la mayor carga de la pierna; la fíbula también transmite parte de la carga y estabiliza el tobillo. El músculo esquelético es estriado y de control somático; el cardíaco es estriado e involuntario, con uno y ocasionalmente dos núcleos centrales; el liso no tiene sarcómeros organizados.
\subsection{Representación conceptual y cuadro de síntesis}
\begin{longtable}{p{0.27\linewidth}p{0.65\linewidth}}
\toprule Concepto & Relación o función\\\midrule
Químico & Moléculas e iones\\
Celular & Células especializadas\\
Tisular & Epitelial, conectivo, muscular y nervioso\\
Órgano & Combinación de tejidos\\
Sistema & Órganos coordinados\\
Organismo & Integración de sistemas\\
\bottomrule
\end{longtable}
\begin{figure}[htbp]
\centering
\begin{tikzpicture}[node distance=5mm, every node/.style={draw,rounded corners=2pt,text width=10cm,align=center,font=\small}]
\node (step0) {Químico: Moléculas e iones};
\node[below=of step0] (step1) {Celular: Células especializadas};
\draw[-{Stealth}] (step0) -- (step1);
\node[below=of step1] (step2) {Tisular: Epitelial, conectivo, muscular y nervioso};
\draw[-{Stealth}] (step1) -- (step2);
\node[below=of step2] (step3) {Órgano: Combinación de tejidos};
\draw[-{Stealth}] (step2) -- (step3);
\node[below=of step3] (step4) {Sistema: Órganos coordinados};
\draw[-{Stealth}] (step3) -- (step4);
\node[below=of step4] (step5) {Organismo: Integración de sistemas};
\draw[-{Stealth}] (step4) -- (step5);
\end{tikzpicture}
\caption{Esquema didáctico original; las flechas indican orden de lectura o relación conceptual, no causalidad universal.}
\end{figure}

\clearpage


\section{Módulo Teórico: Anatomía de Alto Rendimiento}
Este apartado reúne conceptos y relaciones anatómicas para el repaso local. Anatomía es una organización didáctica de este material, no un módulo independiente de EXANI-II ni garantía de cobertura completa de Premedicina.

\subsection{Planimetría y Términos de Orientación}
La anatomía exige un lenguaje universal basado en la \textbf{posición anatómica estándar} (bipedestación, mirada al frente, extremidades superiores pendientes con las palmas hacia adelante).
\begin{itemize}
    \item \textbf{Planos Anatómicos:}
    \begin{itemize}
        \item \textbf{Plano sagital:} Divide en porciones derecha e izquierda; el plano medio las divide en mitades. La flexión y extensión se describen principalmente en este plano alrededor de un eje transversal.
        \item \textbf{Plano coronal (frontal):} Divide en porciones anterior y posterior. La abducción y aducción se describen principalmente en este plano alrededor de un eje anteroposterior.
        \item \textbf{Plano Transversal (Axial):} Eje longitudinal. Divide en porciones superior (craneal) e inferior (caudal). Movimientos de rotación.
    \end{itemize}
    \item \textbf{Términos de Relación:}
    \begin{itemize}
        \item \textit{Proximal / Distal:} Frecuente en las extremidades, tomando como referencia la unión con el tronco (ej. el húmero es proximal al radio).
        \item \textit{Medial / Lateral:} Cercanía a la línea media del cuerpo.
        \item \textit{Ipsilateral / Contralateral:} Del mismo lado o del lado opuesto, respectivamente.
    \end{itemize}
\end{itemize}

\subsection{Cavidades Corporales y Membranas Serosas}
El cuerpo se divide en compartimentos que protegen los órganos.
\begin{itemize}
    \item \textbf{Cavidad Dorsal:} Incluye la cavidad craneal (encéfalo) y el conducto vertebral (médula espinal). Recubierta por las meninges (duramadre, aracnoides, piamadre).
    \item \textbf{Cavidad Ventral:}
    \begin{itemize}
        \item \textit{Cavidad Torácica:} Contiene dos cavidades pleurales (pulmones) y el mediastino (contiene el corazón en la cavidad pericárdica, timo, esófago, tráquea).
        \item \textit{Cavidad Abdominopélvica:} Separada del tórax por el músculo diafragma. La porción abdominal contiene vísceras digestivas; la porción pélvica contiene vejiga, órganos reproductores y recto.
    \end{itemize}
\end{itemize}
\begin{tcolorbox}[colback=blue!5!white,colframe=blue!75!black,title=Nota de estudio]
Las membranas serosas que recubren estas cavidades tienen dos capas: la \textbf{capa parietal} (toca la pared de la cavidad) y la \textbf{capa visceral} (toca y recubre el órgano directamente). Entre ellas hay líquido seroso lubricante. El peritoneo recubre la cavidad abdominal; la pleura, los pulmones; y el pericardio, el corazón.
\end{tcolorbox}

\subsection{Sistema Óseo (Esqueleto Humano)}
El esqueleto humano consta de 206 huesos y tejido cartilaginoso.
\subsubsection{Esqueleto Axial (80 huesos)}
Forma el eje central del cuerpo, diseñado para soporte y protección.
\begin{itemize}
    \item \textbf{Cráneo (22):} 8 craneales (frontal, 2 parietales, 2 temporales, occipital, esfenoides, etmoides) y 14 faciales (maxilares, cigomáticos, nasales, mandíbula, etc.). \textit{Nota: El hueso hioides no se articula con ningún otro hueso y sostiene la lengua.}
    \item \textbf{Columna Vertebral (26):} 7 cervicales (C1=Atlas, C2=Axis), 12 torácicas (articula con costillas), 5 lumbares, 1 sacro (5 vértebras fusionadas), 1 cóccix.
    \item \textbf{Caja Torácica (25):} 1 esternón y 24 costillas (12 pares: 1-7 verdaderas unidas al esternón por cartílago propio; 8-10 falsas unidas al cartílago superior; 11-12 flotantes sin unión anterior).
\end{itemize}

\subsubsection{Esqueleto Apendicular (126 huesos)}
Diseñado para la movilidad (palancas musculares).
\begin{itemize}
    \item \textbf{Cintura Escapular (4):} Clavículas y escápulas.
    \item \textbf{Extremidad Superior (60):} Húmero, Radio (lateral, lado del pulgar), Cúbito/Ulna (medial). Carpo (8 huesos: escafoides, semilunar, piramidal, pisiforme, trapecio, trapezoide, grande, ganchoso). Metacarpos (5), Falanges (14).
    \item \textbf{Cintura Pélvica (2):} Huesos coxales (fusión de ilion, isquion y pubis).
    \item \textbf{Extremidad Inferior (60):} Fémur (hueso más largo/pesado), Rótula/Patela, Tibia (medial y de soporte de peso), Peroné/Fíbula (lateral). Tarso (7 huesos, el calcáneo forma el talón), Metatarsos (5), Falanges (14).
\end{itemize}

\subsection{Sistema Muscular: Mecánica y Estructura}
Los músculos generan fuerza a través del ciclo de los puentes cruzados.
\begin{itemize}
    \item \textbf{Estructura:} El músculo estriado está envuelto por epimisio. Se divide en fascículos (perimisio) y fibras musculares (endomisio).
    \item \textbf{Unidad contráctil:} La \textbf{sarcómera} (limitada por dos discos Z). Contiene miofilamentos gruesos (\textbf{miosina}) y delgados (\textbf{actina}, troponina, tropomiosina).
    \item \textbf{Contracción:} El potencial de acción favorece la liberación de calcio ($\mathrm{Ca}^{2+}$) del retículo sarcoplásmico. El $\mathrm{Ca}^{2+}$ se une a la troponina, moviendo la tropomiosina y exponiendo los sitios de unión en la actina para que la miosina, utilizando ATP, genere la contracción.
\end{itemize}
\begin{tcolorbox}[colback=red!5!white,colframe=red!75!black,title=Clasificación del Tejido Muscular]
1. \textbf{Esquelético:} Estriado, multinucleado periférico, control voluntario somático.\\
2. \textbf{Liso:} No estriado (sin sarcómeras), uninucleado central, control involuntario autónomo (vísceras, vasos sanguíneos).\\
3. \textbf{Cardíaco:} Estriado, habitualmente un núcleo central y ocasionalmente dos, discos intercalares con uniones mecánicas y comunicantes, actividad involuntaria modulada por el sistema autónomo.
\end{tcolorbox}

\newpage
\section{Banco de Reactivos: Actividad 2: banco original de práctica}
Estas preguntas son ejercicios originales de práctica; no son reactivos oficiales ni se ha verificado su equivalencia de dificultad con el EXANI-II. Cada una incluye una respuesta propuesta y su explicación para estudio.

\begin{enumerate}[label=\textbf{\arabic*.}]
    \item \textbf{Un paciente llega a urgencias con una fractura en la diáfisis del hueso largo ubicado en la región lateral del antebrazo (tomando como base la posición anatómica). ¿Qué hueso se ha fracturado?}
    \begin{enumerate}
        \item[A)] Cúbito (Ulna)
        \item[B)] Radio
        \item[C)] Húmero
        \item[D)] Peroné
    \end{enumerate}
    \textit{Respuesta correcta: B. En posición anatómica (palmas hacia adelante), el radio se encuentra en la posición lateral (del lado del pulgar) del antebrazo. El cúbito es medial. El húmero está en el brazo. El peroné es lateral, pero en la pierna.}

    \item \textbf{Durante una intervención quirúrgica abdominal, el cirujano corta primero la capa serosa que recubre directamente la pared interna de la cavidad antes de llegar a los órganos. ¿Qué estructura cortó?}
    \begin{enumerate}
        \item[A)] Peritoneo visceral
        \item[B)] Pleura parietal
        \item[C)] Peritoneo parietal
        \item[D)] Omento mayor
    \end{enumerate}
    \textit{Respuesta correcta: C. La membrana que recubre la cavidad abdominopélvica es el peritoneo. La hoja que recubre la "pared" es la parietal; la que recubre el órgano es la visceral. La pleura pertenece a los pulmones.}

    \item \textbf{¿Cuál es el plano anatómico que permite observar un corte donde se aprecian simultáneamente la región torácica y la región lumbar en un solo lado (perfil) del cuerpo?}
    \begin{enumerate}
        \item[A)] Sagital
        \item[B)] Coronal
        \item[C)] Transversal
        \item[D)] Oblicuo
    \end{enumerate}
    \textit{Respuesta correcta: A. Un corte sagital divide en porciones derecha e izquierda y puede mostrar regiones torácica y lumbar en perfil; solo el corte medio produce mitades iguales. El coronal divide en anterior y posterior y el transversal en superior e inferior.}

    \item \textbf{De los siguientes huesos del cráneo humano, seleccione aquel que sea un hueso impar:}
    \begin{enumerate}
        \item[A)] Parietal
        \item[B)] Temporal
        \item[C)] Esfenoides
        \item[D)] Cigomático
    \end{enumerate}
    \textit{Respuesta correcta: C. El cráneo tiene huesos pares (2 parietales, 2 temporales) e impares (1 frontal, 1 occipital, 1 esfenoides, 1 etmoides). El cigomático es facial y es par (pómulos).}

    \item \textbf{¿Qué describe la unión anterior de las costillas vertebrocondrales 8, 9 y 10, distinguiéndolas de las flotantes 11 y 12?}
    \begin{enumerate}
        \item[A)] No articularse con la columna vertebral torácica.
        \item[B)] No poseer cartílago costal.
        \item[C)] Articularse de manera indirecta con el esternón uniéndose al cartílago de la costilla suprayacente.
        \item[D)] Flotar libremente en la cavidad abdominal.
    \end{enumerate}
    \textit{Respuesta correcta: C. Las costillas 8--10 se unen indirectamente al esternón mediante el arco costal. Las 11--12 no tienen unión anterior y se denominan flotantes; también se incluyen entre las falsas en la clasificación general. Todas se articulan posteriormente con vértebras torácicas.}

    \item \textbf{¿Qué estructura muscular es la principal responsable de la respiración tranquila (basal), separando a su vez la cavidad torácica de la abdominal?}
    \begin{enumerate}
        \item[A)] Músculos intercostales externos
        \item[B)] Diafragma
        \item[C)] Recto abdominal
        \item[D)] Esternocleidomastoideo
    \end{enumerate}
    \textit{Respuesta correcta: B. El diafragma es el músculo primordial de la inspiración y la división anatómica entre tórax y abdomen. Los intercostales apoyan, y el ECM es accesorio (para respiración forzada).}

    \item \textbf{Al observar una muestra de tejido muscular bajo el microscopio, se identifican células cilíndricas largas, con múltiples núcleos localizados en la periferia de la célula y presencia de estrías. Este tejido corresponde a músculo:}
    \begin{enumerate}
        \item[A)] Liso visceral
        \item[B)] Cardíaco
        \item[C)] Esquelético
        \item[D)] Multiunitario
    \end{enumerate}
    \textit{Respuesta correcta: C. El músculo esquelético es el único que es multinucleado y cuyos núcleos están en la periferia, además de ser estriado. El cardíaco es estriado, habitualmente con uno y ocasionalmente dos núcleos centrales. El liso no tiene estrías.}

    \item \textbf{Fisiológicamente, ¿qué ion intracelular es indispensable para que los filamentos de miosina puedan unirse a la actina e iniciar la contracción muscular?}
    \begin{enumerate}
        \item[A)] Sodio ($\mathrm{Na}^{+}$)
        \item[B)] Potasio ($\mathrm{K}^{+}$)
        \item[C)] Calcio ($\mathrm{Ca}^{2+}$)
        \item[D)] Cloruro ($\mathrm{Cl}^{-}$)
    \end{enumerate}
    \textit{Respuesta correcta: C. El Calcio se libera del retículo sarcoplásmico, se une a la troponina y mueve la tropomiosina para despejar los sitios de unión en la actina. El Na+ y K+ participan en el potencial de acción previo.}

    \item \textbf{La unidad funcional básica de contracción en el músculo estriado, delimitada por dos líneas Z sucesivas, se denomina:}
    \begin{enumerate}
        \item[A)] Sarcoplasma
        \item[B)] Fascículo
        \item[C)] Sarcolema
        \item[D)] Sarcómera
    \end{enumerate}
    \textit{Respuesta correcta: D. La sarcómera es la unidad funcional mínima. El sarcolema es la membrana celular, el sarcoplasma es el citoplasma y el fascículo es un paquete de fibras.}

    \item \textbf{¿Cuál de las siguientes afirmaciones describe correctamente la primera vértebra cervical (C1 o Atlas)?}
    \begin{enumerate}
        \item[A)] Posee una apófisis odontoides prominente.
        \item[B)] Carece de cuerpo vertebral y de apófisis espinosa.
        \item[C)] Es la vértebra con el cuerpo más grande para soportar el peso.
        \item[D)] Se articula con las costillas superiores.
    \end{enumerate}
    \textit{Respuesta correcta: B. El Atlas (C1) es único por ser un anillo óseo sin cuerpo ni apófisis espinosa. La apófisis odontoides pertenece al Axis (C2). Los cuerpos grandes están en las lumbares.}

    \item \textbf{El mediastino es el compartimento central de la cavidad torácica. ¿Cuál de los siguientes órganos NO se localiza en el mediastino?}
    \begin{enumerate}
        \item[A)] Tráquea
        \item[B)] Pulmones
        \item[C)] Corazón
        \item[D)] Esófago
    \end{enumerate}
    \textit{Respuesta correcta: B. Los pulmones no están en el mediastino; ocupan sus propias cavidades (cavidades pleurales derecha e izquierda) a los lados del mediastino.}

    \item \textbf{Se le llama "hueso sesamoideo" a aquel que se desarrolla dentro de un tendón sometido a fricción y tensión. ¿Cuál es el hueso sesamoideo más grande del cuerpo humano?}
    \begin{enumerate}
        \item[A)] Pisiforme
        \item[B)] Astrágalo
        \item[C)] Rótula (Patela)
        \item[D)] Hioides
    \end{enumerate}
    \textit{Respuesta correcta: C. La rótula se forma dentro del tendón del músculo cuádriceps femoral, protegiendo la articulación de la rodilla y mejorando el apalancamiento mecánico.}

    \item \textbf{La articulación de la cadera (coxofemoral) y la articulación del hombro (glenohumeral) son ejemplos típicos de articulaciones sinoviales de tipo:}
    \begin{enumerate}
        \item[A)] Bisagra (Tróclea)
        \item[B)] Esférica (Enartrosis)
        \item[C)] Pivote (Trocoide)
        \item[D)] Plana (Artrodia)
    \end{enumerate}
    \textit{Respuesta correcta: B. Las enartrosis o articulaciones esféricas permiten el mayor grado de libertad (movimiento en todos los ejes y planos). La bisagra es el codo; el pivote es la atlantoaxoidea.}

    \item \textbf{A nivel del carpo (muñeca), ¿qué par de huesos se articulan directamente con el hueso radio para formar la articulación radiocarpiana?}
    \begin{enumerate}
        \item[A)] Escafoides y semilunar
        \item[B)] Trapecio y trapezoide
        \item[C)] Grande y ganchoso
        \item[D)] Piramidal y pisiforme
    \end{enumerate}
    \textit{Respuesta correcta: A. La superficie distal del radio contacta directamente con el escafoides y el semilunar. El piramidal se relaciona con el disco articular, no directamente con el radio. Las otras parejas propuestas no contactan directamente con el radio.}

    \item \textbf{En el desarrollo embrionario y el crecimiento óseo, ¿cómo se le conoce a la zona de cartílago hialino ubicada entre la diáfisis y la epífisis de un hueso largo, responsable del crecimiento en longitud?}
    \begin{enumerate}
        \item[A)] Periostio
        \item[B)] Endostio
        \item[C)] Placa epifisaria (Fisis)
        \item[D)] Médula amarilla
    \end{enumerate}
    \textit{Respuesta correcta: C. La placa o cartílago epifisario (fisis) permite el crecimiento longitudinal. Cuando se osifica al final de la pubertad, se le llama línea epifisaria y el hueso deja de crecer en longitud.}

    \item \textbf{¿Qué estructura ósea forma la porción posterior e inferior del hueso coxal (la parte sobre la cual nos sentamos)?}
    \begin{enumerate}
        \item[A)] Ilion
        \item[B)] Isquion
        \item[C)] Pubis
        \item[D)] Sacro
    \end{enumerate}
    \textit{Respuesta correcta: B. El hueso coxal consta de ilion (superior), pubis (anteroinferior) e isquion (posteroinferior). La tuberosidad isquiática es la estructura ósea que soporta el peso al sentarse.}

    \item \textbf{Los discos intercalares son estructuras microscópicas ricas en desmosomas y uniones comunicantes (gap). ¿En qué tejido se encuentran de manera exclusiva?}
    \begin{enumerate}
        \item[A)] Tejido óseo compacto
        \item[B)] Músculo cardíaco
        \item[C)] Músculo liso
        \item[D)] Cartílago fibroso
    \end{enumerate}
    \textit{Respuesta correcta: B. Los discos intercalares unen mecánicamente y eléctricamente a las fibras musculares cardíacas, permitiendo que el miocardio se contraiga como un sincitio funcional.}

    \item \textbf{Al realizar una abducción del brazo, la extremidad se mueve:}
    \begin{enumerate}
        \item[A)] Alejándose de la línea media del cuerpo.
        \item[B)] Acercándose hacia la línea media del cuerpo.
        \item[C)] Flexionándose hacia adelante.
        \item[D)] Rotando sobre su propio eje.
    \end{enumerate}
    \textit{Respuesta correcta: A. Abducción significa "alejar" del plano medio sagital. Aducción significa "acercar". Flexión reduce el ángulo entre los huesos.}

    \item \textbf{¿Qué hueso del lado lateral de la pierna forma el maléolo lateral, se articula con tibia y talo y no directamente con el fémur?}
    \begin{enumerate}
        \item[A)] Tibia
        \item[B)] Fémur
        \item[C)] Peroné (Fíbula)
        \item[D)] Calcáneo
    \end{enumerate}
    \textit{Respuesta correcta: C. La fíbula o peroné forma el maléolo lateral y contacta con tibia y talo, no con el fémur. La tibia es medial; el fémur corresponde al muslo y el calcáneo al pie. La pregunta no requiere asignar un porcentaje de carga a la fíbula.}

    \item \textbf{La capa de tejido conectivo que recubre a la fibra muscular de forma individual es el:}
    \begin{enumerate}
        \item[A)] Epimisio
        \item[B)] Perimisio
        \item[C)] Endomisio
        \item[D)] Sarcolema
    \end{enumerate}
    \textit{Respuesta correcta: C. El Endomisio recubre la fibra muscular individual. El perimisio recubre los fascículos (grupos de fibras) y el epimisio recubre todo el músculo. El sarcolema es la membrana plasmática, no tejido conectivo externo.}

\end{enumerate}


\section{Conclusiones}
El material permite repasar conceptos y contrastar respuestas. La clave no acredita desempeño del aspirante ni aprobación institucional. Deben revisarse la convocatoria y el temario aplicables antes de usarlo como plan definitivo de ingreso.
\clearpage\section{Fuentes de apoyo y alcance de consulta}
OpenStax. Anatomy and Physiology 2e: 1.6 Anatomical Terminology. Consulta: 10 de octubre de 2026. \url{https://openstax.org/books/anatomy-and-physiology-2e/pages/1-6-anatomical-terminology}\par Sección consultada: posición, términos, planos y cavidades. No valida todos los reactivos musculoesqueléticos.\par\medskip
Los originales bibliográficos disponibles se identifican por nombre de archivo y SHA256 en el inventario local. No se atribuyen ediciones, páginas ni lecturas no comprobadas.
\end{document}
